\documentclass[12pt,oneside]{book}
\usepackage[spanish]{babel}
\usepackage{graphicx}
\usepackage{amsmath}
\usepackage{array}
\usepackage{booktabs}
\usepackage[a4paper,margin=2.5cm]{geometry}
\usepackage{etoolbox}
\usepackage{fancyhdr}
\usepackage{lastpage}
\usepackage{xcolor}
\graphicspath{{imagenes/}{Imagenes/}{outputs/}}
\makeatletter
\patchcmd{\@makechapterhead}{\vspace*{50\p@}}{\vspace*{0\p@}}{}{}
\patchcmd{\@makeschapterhead}{\vspace*{50\p@}}{\vspace*{0\p@}}{}{}
\makeatother
\pagestyle{fancy}
\fancyhf{}
\setlength{\headheight}{15pt}
\fancyhead[C]{\textcolor{gray}{\small Práctica 4 - IAA - EPSC}}
\renewcommand{\headrulewidth}{0pt}
\addtolength{\footskip}{8pt}
\fancyfoot[C]{\hspace*{-1.6cm}\textcolor{gray}{Juan Luis Prieto Panadero y Alberto Álvarez Mejías}}
\fancyfoot[R]{\textcolor{gray}{Pagina \thepage\ de \pageref{LastPage}}}
\fancypagestyle{plain}{%
    \fancyhf{}
    \fancyhead[C]{\textcolor{gray}{\small Práctica 4 - Escuela Politécnica Superior de Córdoba}}
    \renewcommand{\headrulewidth}{0pt}
    \fancyfoot[C]{\hspace*{-1.6cm}\textcolor{gray}{Juan Luis Prieto Panadero y Alberto Álvarez Mejías}}
    \fancyfoot[R]{\textcolor{gray}{Pagina \thepage\ de \pageref{LastPage}}}
}

\begin{document}

\begin{titlepage}
\thispagestyle{empty}
\centering

\vspace*{0.5cm}
\noindent
\begin{minipage}[c]{0.45\textwidth}
    \raggedright
    
\includegraphics[height=2.2cm]{logo_uco.png}
\end{minipage}
\begin{minipage}[c]{0.45\textwidth}
    \raggedleft
    
\includegraphics[height=2.2cm]{logotipo-EPSC.png}
\end{minipage}

\vspace{0.9cm}
\rule{0.9\textwidth}{0.4pt}

\vspace{0.9cm}
{\LARGE\bfseries Práctica 4\par}
\vspace{0.25cm}
{\Large\bfseries Introducción al Aprendizaje Automático\par}

\vspace{0.9cm}
\rule{0.9\textwidth}{0.4pt}

\vspace{0.9cm}
{\large\bfseries Escuela Politécnica Superior de Córdoba\par}
\vspace{0.25cm}
{\large\bfseries Universidad de Córdoba\par}
\vspace{0.25cm}
{\large\bfseries Grado de Ingeniería Informática\par}

\vfill
{\large\bfseries Autores:\par}
\vspace{0.2cm}
{\normalsize Juan Luis Prieto Panadero\par}
{\normalsize Alberto Álvarez Mejías\par}

\vspace*{0.5cm}
\end{titlepage}

\clearpage
\setcounter{page}{1}

\counterwithout{section}{chapter}

% =========================================================
% CAPÍTULO / SECCIÓN DE MEMORIA — PRÁCTICA 4
% Clasificación y el Valor del Error (Matriz de Costes)
% =========================================================

\section{Práctica 4: Clasificación y el Valor del Error}

\subsection{Objetivo de la práctica}
En esta práctica se entrena un clasificador probabilístico sencillo para un problema de detección de fraude y se analiza cómo el umbral de decisión afecta a los errores y al coste total. El objetivo no es únicamente medir el rendimiento predictivo del modelo, sino estudiar el impacto real de sus decisiones a partir de una matriz de costes.

\subsection{Descripción del problema}
Se trabaja con un problema de clasificación binaria:
\begin{itemize}
    \item \textbf{Clase 0}: transacción normal.
    \item \textbf{Clase 1}: transacción fraudulenta.
\end{itemize}

En este contexto, los errores no tienen el mismo coste:
\begin{itemize}
    \item \textbf{Falso Negativo (FN)}: no detectar un fraude \(\rightarrow 200\,€\)
    \item \textbf{Falso Positivo (FP)}: bloquear una transacción legítima \(\rightarrow 10\,€\)
\end{itemize}

Por tanto, la decisión final no debe evaluarse solo por el número de aciertos, sino por el coste económico asociado a los errores.

% ---------------------------------------------------------
\subsection{Tarea 1: Análisis con el umbral por defecto (\(\tau = 0.5\))}
% ---------------------------------------------------------

\subsubsection{Modelo utilizado}
En esta práctica se ha utilizado una \textbf{Regresión Logística} como clasificador base. El modelo se ha entrenado con el conjunto de entrenamiento y se ha evaluado sobre el conjunto de prueba.

\subsubsection{Matriz de confusión}
A continuación se muestra la matriz de confusión obtenida usando el umbral estándar \(\tau = 0.5\).

\begin{figure}[H]
    \centering
    \includegraphics[width=0.48\textwidth]{img/matriz_confusion_umbral_0_5.png}
    \caption{Matriz de confusión para el umbral \(\tau = 0.5\).}
    \label{fig:mc_05}
\end{figure}

A partir de esta matriz, se obtienen los siguientes valores:
\begin{itemize}
    \item Verdaderos Positivos (TP): \texttt{[SUSTITUIR]}
    \item Verdaderos Negativos (TN): \texttt{[SUSTITUIR]}
    \item Falsos Positivos (FP): \texttt{[SUSTITUIR]}
    \item Falsos Negativos (FN): \texttt{[SUSTITUIR]}
\end{itemize}

\subsubsection{Cálculo del coste total}
El coste total se calcula mediante la expresión:
\[
\text{Coste Total} = (FP \cdot 10) + (FN \cdot 200)
\]

Sustituyendo los valores obtenidos:
\[
\text{Coste Total} = (\texttt{[FP]} \cdot 10) + (\texttt{[FN]} \cdot 200) = \texttt{[SUSTITUIR]} \,€
\]

\subsubsection{Interpretación}
\textbf{Comentario del alumno:}  
\textit{Explica aquí si el umbral por defecto de 0.5 parece razonable o no en este contexto. No te limites a decir si el modelo acierta mucho o poco; comenta especialmente si está controlando adecuadamente los errores más costosos.}

\vspace{0.3cm}
\noindent\textbf{Guía para redactar:}
\begin{itemize}
    \item ¿El modelo comete muchos falsos negativos?
    \item ¿El coste total parece alto o bajo para este problema?
    \item ¿Usarías este umbral en un entorno real sin revisarlo?
\end{itemize}

% ---------------------------------------------------------
\subsection{Tarea 2: Estudio de probabilidades y umbrales}
% ---------------------------------------------------------

\subsubsection{Análisis de distintos umbrales}
En lugar de trabajar únicamente con la clase predicha, se han utilizado las probabilidades estimadas por el modelo para la clase positiva y se han evaluado distintos valores del umbral \(\tau\) entre 0 y 1, en pasos de 0.05.

La Tabla~\ref{tab:umbrales_coste} resume los resultados obtenidos para varios umbrales.

\begin{table}[H]
\centering
\caption{Resumen de resultados para distintos umbrales.}
\label{tab:umbrales_coste}
\begin{tabular}{c c c c}
\toprule
\textbf{Umbral} & \textbf{FP} & \textbf{FN} & \textbf{Coste Total (€)} \\
\midrule
0.00 & \texttt{[ ]} & \texttt{[ ]} & \texttt{[ ]} \\
0.05 & \texttt{[ ]} & \texttt{[ ]} & \texttt{[ ]} \\
0.10 & \texttt{[ ]} & \texttt{[ ]} & \texttt{[ ]} \\
0.15 & \texttt{[ ]} & \texttt{[ ]} & \texttt{[ ]} \\
0.20 & \texttt{[ ]} & \texttt{[ ]} & \texttt{[ ]} \\
\vdots & \vdots & \vdots & \vdots \\
1.00 & \texttt{[ ]} & \texttt{[ ]} & \texttt{[ ]} \\
\bottomrule
\end{tabular}
\end{table}

\subsubsection{Observaciones}
\textbf{Comentario del alumno:}  
\textit{Describe aquí cómo cambian los falsos positivos y los falsos negativos al mover el umbral.}

\vspace{0.3cm}
\noindent\textbf{Guía para redactar:}
\begin{itemize}
    \item ¿Qué ocurre cuando el umbral baja?
    \item ¿Qué ocurre cuando el umbral sube?
    \item ¿Qué compromiso aparece entre FP y FN?
\end{itemize}

% ---------------------------------------------------------
\subsection{Tarea 3: Optimización del coste}
% ---------------------------------------------------------

\subsubsection{Gráfica de coste frente a umbral}
La Figura~\ref{fig:coste_umbral} muestra la evolución del coste total en función del umbral de decisión.

\begin{figure}[H]
    \centering
    \includegraphics[width=0.72\textwidth]{img/coste_vs_umbral.png}
    \caption{Coste total en función del umbral de decisión.}
    \label{fig:coste_umbral}
\end{figure}

\subsubsection{Umbral óptimo}
A partir de la gráfica y de la tabla anterior, el umbral que minimiza el coste total es:

\[
\tau^* = \texttt{[SUSTITUIR]}
\]

El coste total mínimo asociado es:

\[
\text{Coste mínimo} = \texttt{[SUSTITUIR]} \,€
\]

\subsubsection{Interpretación}
\textbf{Comentario del alumno:}  
\textit{Explica aquí si el umbral óptimo es mayor o menor que 0.5 y por qué ese valor tiene sentido en este problema.}

\vspace{0.3cm}
\noindent\textbf{Guía para redactar:}
\begin{itemize}
    \item ¿Por qué el mínimo no tiene por qué aparecer en 0.5?
    \item ¿Qué relación tiene ese valor con la diferencia de costes entre FP y FN?
    \item ¿Qué estrategia parece más razonable para el banco: conservadora o permisiva?
\end{itemize}

% ---------------------------------------------------------
\subsection{Tarea 4: Curvas ROC y Precision-Recall}
% ---------------------------------------------------------

\subsubsection{Curva ROC}
La Figura~\ref{fig:roc} muestra la curva ROC obtenida para el modelo.

\begin{figure}[H]
    \centering
    \includegraphics[width=0.58\textwidth]{img/curva_roc.png}
    \caption{Curva ROC del clasificador.}
    \label{fig:roc}
\end{figure}

\subsubsection{Curva Precision-Recall}
La Figura~\ref{fig:pr} muestra la curva Precision-Recall del modelo.

\begin{figure}[H]
    \centering
    \includegraphics[width=0.58\textwidth]{img/curva_precision_recall.png}
    \caption{Curva Precision-Recall del clasificador.}
    \label{fig:pr}
\end{figure}

\subsubsection{Discusión}
\textbf{Comentario del alumno:}  
\textit{Si el dataset está desbalanceado, razona cuál de las dos curvas resulta más útil para valorar la utilidad real del modelo en este problema.}

\vspace{0.3cm}
\noindent\textbf{Guía para redactar:}
\begin{itemize}
    \item ¿Qué información aporta la curva ROC?
    \item ¿Qué información aporta la curva Precision-Recall?
    \item ¿Cuál refleja mejor el comportamiento sobre la clase minoritaria?
\end{itemize}

% ---------------------------------------------------------
\subsection{Decisión final de ingeniería}
% ---------------------------------------------------------

Imaginando que formo parte del equipo técnico del banco, la política de decisión que recomendaría sería la siguiente:

\begin{quote}
\texttt{[REDACTAR AQUÍ UNA PROPUESTA BREVE Y ARGUMENTADA]}
\end{quote}

\noindent Esta decisión se basa en:
\begin{itemize}
    \item el coste asociado a los falsos positivos y falsos negativos;
    \item el umbral que minimiza el coste total;
    \item el equilibrio entre seguridad y molestias para el cliente;
    \item la interpretación conjunta de las curvas ROC y Precision-Recall.
\end{itemize}

% ---------------------------------------------------------
\subsection{Conclusiones}
% ---------------------------------------------------------

\textbf{Comentario final del alumno:}  
\textit{Resume aquí qué has aprendido con esta práctica y explica por qué en clasificación no todos los errores valen lo mismo.}

\vspace{0.3cm}
\noindent\textbf{Aspectos que conviene mencionar:}
\begin{itemize}
    \item que el umbral de decisión no es fijo ni universal;
    \item que cambiar el umbral modifica la relación entre FP y FN;
    \item que el contexto del problema debe influir en la evaluación;
    \item que una métrica general no siempre refleja el impacto real de un modelo.
\end{itemize} 

\end{document}